% GENERATED FILE. Edit canonical lessons, then run npm run generate:books.
% canonical-source-hash: fnv1a64:284a278256744d10
% canonical-lessons: GE-C347-zufall, GE-C347-zufaellig, GE-C347-verhindern, GE-C347-vermeiden, GE-C347-weshalb

\chapter{Zufall oder nicht?}
\label{ch:ge-zufall-oder-nicht}
\hlchaptermodality{347}

\begin{chapteropening}
\textbf{By the end of this chapter:} \emph{I can say that something happened by chance, ask why with weshalb, and say how to prevent or avoid a problem.}

Complete the last lesson of chapter 347: I can say that something happened by chance, ask why with weshalb, and say how to prevent or avoid a problem.
\end{chapteropening}

\section[der Zufall]{der Zufall — the chance, the coincidence}

\label{lesson:GE-C347-zufall}



\begin{quote}
Before the new one: say the German for on purpose, then the German for the occasion.
\end{quote}

\subsection*{You'll want to know: der Zufall}
\textbf{der Zufall} — \textquotedblleft{}the chance, the coincidence\textquotedblright{}.

\begin{quote}
Was für ein Zufall! Wir sind im selben Hotel.
\end{quote}

\begin{quote}
\emph{What a coincidence! We are in the same hotel.}
\end{quote}

Plural: \emph{die Zufälle}. \emph{Durch Zufall} — by chance.

\begin{cousinweb}
\emph{Zufall} is built from \emph{zu} \textquotedblleft{}to\textquotedblright{} and \emph{Fall} \textquotedblleft{}fall, case\textquotedblright{}, from Old High German \emph{val} (EtymDB): what falls to you. It is usually explained as a rendering of Latin \emph{accidēns}, from \emph{accidere} \textquotedblleft{}to fall upon, to happen\textquotedblright{} (Lewis \& Short).
\end{cousinweb}

\begin{grammarlens}[title={What you've built}]
One more word for telling, explaining and describing at B1.
\end{grammarlens}

\subsection*{Your turn}
\begin{itemize}
\raggedright
  \item \emph{Say it:} \emph{der Zufall}
  \item \emph{Say it:} \emph{der Zufall}, once more
  \item \emph{Say it:} a sentence of your own with \emph{der Zufall}
  \item \emph{Recall:} say the German for on purpose, then the German for the occasion, then say \emph{der Zufall} again
\end{itemize}

\begin{tcolorbox}[breakable,colback=teal!4,colframe=teal!35!black,title={Before you move on}]
What does der Zufall mean? (\textquotedblleft{}The chance, the coincidence\textquotedblright{}.) Say it once more.
\end{tcolorbox}
\section[zufällig]{zufällig — by chance, accidental}

\label{lesson:GE-C347-zufaellig}



\begin{quote}
Before the new one: say the German for the occasion, then the German for the chance.
\end{quote}

\subsection*{You'll want to know: zufällig}
\textbf{zufällig} — \textquotedblleft{}by chance, accidental\textquotedblright{}.

\begin{quote}
Ich habe meine alte Lehrerin zufällig im Supermarkt getroffen.
\end{quote}

\begin{quote}
\emph{I met my old teacher by chance in the supermarket.}
\end{quote}

An adverb and an adjective: \emph{ein zufälliges Treffen} — a chance meeting.

\begin{cousinweb}
\emph{zufällig} is built with \emph{-ig} on \emph{Zufall} \textquotedblleft{}chance\textquotedblright{}, from \emph{zu} and \emph{Fall} \textquotedblleft{}fall\textquotedblright{}, Old High German \emph{val} (EtymDB).
\end{cousinweb}

\begin{grammarlens}[title={What you've built}]
One more word for telling, explaining and describing at B1.
\end{grammarlens}

\subsection*{Your turn}
\begin{itemize}
\raggedright
  \item \emph{Say it:} \emph{zufällig}
  \item \emph{Say it:} \emph{zufällig}, once more
  \item \emph{Say it:} a sentence of your own with \emph{zufällig}
  \item \emph{Recall:} say the German for the occasion, then the German for the chance, then say \emph{zufällig} again
\end{itemize}

\begin{tcolorbox}[breakable,colback=teal!4,colframe=teal!35!black,title={Before you move on}]
What does zufällig mean? (\textquotedblleft{}By chance, accidental\textquotedblright{}.) Say it once more.
\end{tcolorbox}
\section[verhindern]{verhindern — to prevent}

\label{lesson:GE-C347-verhindern}



\begin{quote}
Before the new one: say the German for the chance, then the German for by chance.
\end{quote}

\subsection*{You'll want to know: verhindern}
\textbf{verhindern} — \textquotedblleft{}to prevent\textquotedblright{}.

\begin{quote}
Der Regen hat das Spiel nicht verhindert.
\end{quote}

\begin{quote}
\emph{The rain did not prevent the game.}
\end{quote}

Inseparable: \emph{er verhindert, er hat verhindert}. \emph{Verhindert sein} — to be unable to come.

\begin{cousinweb}
\emph{verhindern} is built with \emph{ver-} on \emph{hindern} \textquotedblleft{}to hinder\textquotedblright{}, a verb from the family of \emph{hinter} \textquotedblleft{}behind\textquotedblright{}: you hold a thing back. EtymDB has no history for the verbs; Proto-Germanic \emph{hinder} \textquotedblleft{}behind\textquotedblright{} is the usual source.
\end{cousinweb}

\begin{grammarlens}[title={What you've built}]
One more word for telling, explaining and describing at B1.
\end{grammarlens}

\subsection*{Your turn}
\begin{itemize}
\raggedright
  \item \emph{Say it:} \emph{verhindern}
  \item \emph{Say it:} \emph{verhindern}, once more
  \item \emph{Say it:} a sentence of your own with \emph{verhindern}
  \item \emph{Recall:} say the German for the chance, then the German for by chance, then say \emph{verhindern} again
\end{itemize}

\begin{tcolorbox}[breakable,colback=teal!4,colframe=teal!35!black,title={Before you move on}]
What does verhindern mean? (\textquotedblleft{}To prevent\textquotedblright{}.) Say it once more.
\end{tcolorbox}
\section[vermeiden]{vermeiden — to avoid}

\label{lesson:GE-C347-vermeiden}



\begin{quote}
Before the new one: say the German for by chance, then the German for to prevent.
\end{quote}

\subsection*{You'll want to know: vermeiden}
\textbf{vermeiden} — \textquotedblleft{}to avoid\textquotedblright{}.

\begin{quote}
Ich fahre früh los, um den Verkehr zu vermeiden.
\end{quote}

\begin{quote}
\emph{I set off early to avoid the traffic.}
\end{quote}

Inseparable and strong: \emph{er vermeidet, er hat vermieden}.

\begin{cousinweb}
\emph{vermeiden} is built with \emph{ver-} on \emph{meiden} \textquotedblleft{}to keep away from\textquotedblright{}, inherited from Old High German \emph{mīdan} and Proto-Germanic \emph{mīþaną} \textquotedblleft{}to avoid\textquotedblright{} (EtymDB).
\end{cousinweb}

\begin{grammarlens}[title={What you've built}]
One more word for telling, explaining and describing at B1.
\end{grammarlens}

\subsection*{Your turn}
\begin{itemize}
\raggedright
  \item \emph{Say it:} \emph{vermeiden}
  \item \emph{Say it:} \emph{vermeiden}, once more
  \item \emph{Say it:} a sentence of your own with \emph{vermeiden}
  \item \emph{Recall:} say the German for by chance, then the German for to prevent, then say \emph{vermeiden} again
\end{itemize}

\begin{tcolorbox}[breakable,colback=teal!4,colframe=teal!35!black,title={Before you move on}]
What does vermeiden mean? (\textquotedblleft{}To avoid\textquotedblright{}.) Say it once more.
\end{tcolorbox}
\section[weshalb]{weshalb — why, for what reason}

\label{lesson:GE-C347-weshalb}



\begin{quote}
Before the new one: say the German for to prevent, then the German for to avoid.
\end{quote}

\subsection*{You'll want to know: weshalb}
\textbf{weshalb} — \textquotedblleft{}why, for what reason\textquotedblright{}.

\begin{quote}
Niemand weiß, weshalb der Laden heute geschlossen ist.
\end{quote}

\begin{quote}
\emph{Nobody knows why the shop is closed today.}
\end{quote}

A question word like \emph{warum}. In an indirect question it sends the verb to the end: \emph{…, weshalb er geschlossen ist}.

\begin{cousinweb}
\emph{weshalb} joins \emph{wes}, the old genitive of \emph{was} \textquotedblleft{}what\textquotedblright{}, to \emph{halb} in its old sense \textquotedblleft{}on the side of, because of\textquotedblright{}: \textquotedblleft{}because of what\textquotedblright{}. EtymDB has no entry, so this follows the usual account.
\end{cousinweb}

\begin{grammarlens}[title={What you've built}]
One more word for telling, explaining and describing at B1.
\end{grammarlens}

\subsection*{Your turn}
\begin{itemize}
\raggedright
  \item \emph{Say it:} \emph{weshalb}
  \item \emph{Say it:} \emph{weshalb}, once more
  \item \emph{Say it:} a sentence of your own with \emph{weshalb}
  \item \emph{Recall:} say the German for to prevent, then the German for to avoid, then say \emph{weshalb} again
\end{itemize}

\begin{tcolorbox}[breakable,colback=teal!4,colframe=teal!35!black,title={Before you move on}]
What does weshalb mean? (\textquotedblleft{}Why, for what reason\textquotedblright{}.) Say it once more.
\end{tcolorbox}
